\documentclass[]{../../../../tex/classes/styledArticle}
\usepackage{../../../../tex/packages/hebrewSupport}

\author{שחר פרץ}
\title{\textit{היסטוריה 48}}
\begin{document}
	\maketitle
	בשיעור הקרוב נעבור על תיקונים מהספר בהתבסס על מחקרים חדשים ופרוטוקולים שיצאו. הספר וצורת הלימוד עד היום ''לא נכונה במקרה הטוב``. 
	
	\subsubsection*{תוצאות}
	זכרו – תוצאה, היא מיידית. מצלמים תמונת מצב כפי שהייתה ביום סיום המלחמה. בד''כ מדברים על דברים קונקרטיים ולא על תהליכים, כי תהליכים זה השפעות. 
	
	מצד אחד, ישנו נצחון צבאי מוחלט. אבדות כבדות במטוסים ובטנקים, ופגיעה בכושר צה''ל. דיברנו על זה בשיעור הקודם. 
	
	\subsubsection*{השפעות}
	תחושת הבטחון והאמונה המלא בתפיסה הבטחונית נשברה. יש ירידה משמעותית באמונה, והביקורת המשמעותית שמגיעה מלמטה הפכה להיות משמעותית הרבה יותר ממלחמת יום הכיפורים. צריך להבין, שבאותה התקופה ב''מעריב לנוער`` תמיד היו פוסטרים זמרים ולהקות, ובין ששת הימים ויום כיפור היו פוסטרים של קציני צבא כמו רבין. הייתה ממש הערצה אליהם, כמו אל כוכבי תרבות אחרים. השבר הזה משמעותי במיוחד, לעומת מה שהיה לפני כן. 
	
	הוקמה ועדת חקירה ממלכתית, ועדת אגרנט. חקרה את מחדלי המלחמה ומסקנותיה הופנו כנגד הדרג הצבאי. הרמתכ''ל רא''ל אלעזר נאלץ להתפטר ושורה של קצינים בכירים באמ''ן הודחו. מבחינה ציבורית, הציבור לא ממש קיבל את ההמלצות הללו – הוא לא היה מוכן לכך שפוטרים את הדרג המדיני מאחריות. התפתחו מחאות רבות. מוטי אשכנזי, סרן בתפקד על מעוז ''בודפשט`` בצפון התעלה, הוביל לרצף מחאות בארץ בנוגע להתפטרות הסגל המדיני. 
	
	ב־1977 יש מהפך ונוצרו שינויים במפה הפוליטית. זכרו, ב־1964 היו בכירות והמערך על שוב לשלטון, לאחר מכן רבין. לקח מספר שנים עד שהמהפך התרחש. ב־77 עלתה מלחמת הליכוד לשלטון (בין היתר, היחס לפנתרים השחורים. גולדה אמרה עליהם ''הם לא נחמדים`` (שרית אומרת במבטא של גולדה)). 
	
	יש התמודדות קשה על השכול. אלפי הרוגים. הוקדם מרכז מיוחד לאיתור שבויים ונעדרים ונעשו מאמצים רבים בשנים הבאות להתחקות על גורלם [עצה של שרית: לא לכתוב את זה בבחינות. זה די מיידי וקצת עלול להכנס לתוצאות]. היה תהליך של חזרה בתשובה לאחר המלחמה, בדיוק כמו שקורה עכשיו. 
	
	התרחשה הגברת קיטוב בחברה הישראלית בענייני גוף ובטחון. מצד אחד היו אלו שתמכו בהגברת אחיזה של ישראל בשטחי א''י השלמה, לעומת אלו שתמכו בשטחים תמורת שלום. 
	
	ישראל הפכה להיות תלויה בארה''ב מבחינה כלכלית־צבאית. ארה''ב במהלך המלחמה העבירה רכבת אווירית של חימוש, נשק ואף מטוסים. מאותו הרגע ועד היום ישראל תלויה בארה''ב כלכלית. מערבותה של ארה''ב שגברה באותה העת במדינות ערב אפשרה לה לתווך בין ישראל למצריים (ובהסכמים נוספים לאחר מכן). ישראל הפכה לתלויה כלכלית בארה''ב גם בגלל הנזק הכלכלי הכבד שנגרם כאן. 
	
	המלחמה סללה את הדרך לסאדאת להגיע לארץ ולהביא להסכם שלום (הוא השיג בסוף את מה שהוא רצה). 
	
	\subsection*{יואב גלבר}
	יואב גלבר, חבר בועדת אגרנאט שלחם במלחמת יום הכיפורים הוציא ספר על המלחמה. יש מצגת למורים ששרית מציגה עכשיו. 
	
	עד לא מזמן, מהרגע שהסתיימה המלחמה, היו מספר טענות כנגד הדרג המדיני: 
	\begin{itemize}
		\item המלחמה הייתה יכולה להמנע אלמלא גולדה מאיר ומשה דיין היו צרי אופקים ומוגבלים מבחינת ראייתם הכוללת
		\item ממשלת ישראל לא נקטה כל יוזמה כדי למנוע את המלחמה באמצעות הסדרים עם מצריים
		\item הפערים בין ישראל למצריים היו ניתנים לגישור אלמלא עמדתה הנוקשה של ישראל
	\end{itemize}
	כך ''למדנו ולימדנו`` – מצריים יזמה מגעים לתהליך שלום, שמדינת ישראל סירבה לו. ע''פ גלבר, שמתבסס על מחקרים חדשים יותר:
	
	\begin{itemize}
		\item ישראל ניהלה דיפלומטיה עם גורמים רבים: פלשתינים, חוסיין, ושרי חוץ אירופאיים
		\item המצרים נהלו מו''מ \textit{רק עם ארה''ב}, והדברים האלו לא הבשילו מספיק כדי להגיע דרך ארה''ב לגולדה מאיר, כך שהיא לא יכלה לדחות דברים שלא מגיעים אליה
		\item היו פערים משמעותיים בין העמדות של הישראלים למצרים, כלומר הפערים לא היו ניתנים לגישור
	\end{itemize}
	
	מה היו הפערים? מצד מצריים:
	\begin{itemize}
		\item מצרים דרשה הסכם כולל ביחד עם סוריה וירדן. הסכם אזורי, לא רק מדינה אחת. 
		\item התנגדות להסכם עם ישראל כל עוד צה''ל נמצא על אדמות ערביות
		\item פתרון לבעיה הפלסטינית
		\item לא רצו הסכם שלום, אלא הסכם להפסקת לוחמה
		\item התנגדות להכרה במדינה
		\item נסיגה מוחלטת מהשטחים שנכבשו ש־66, ושסיני לא תהיה מפורזת/רצועה מפורזת צרה והדדית
	\end{itemize}
	מצד ישראל: 
	\begin{itemize}
		\item מו''מ ישיר
		\item התקדמות בשלבים
		\item לא לדון בפתרון הבעיה הפלסטינית
		\item ישראל רצתה גבולות בטחון בתוך סיני
		\item פירוז סיני מכוחות מצרים
	\end{itemize}
	
	בספר כתוב ''גישושי שלום נתקלו בסירוב מן הצד הישראלי, שלא ייחס להם כוונות רציניות`` [עמוד 185] וכן ''סאדאת הודיע שהיה מוכן לשאת ולתת על הסכם שלום עם ישראל ואף להסדר ביניים (נסיגה חלקית) שבסופו ישראל תיסוג מכל השטחים שכבשה`` [עמוד 245]. המשפטים האלו \textit{לא נכונים}. הגורמים למלחמה שמצויים בספר הם לא הגורמים. הזלזול בכוח הערבי, פירוש שגוי של המידע המודיעיני, וכו', אינם \textit{גורמים למלחמה}, אלא \textit{הגורמים להפתעה}. יצאה הנחיה מפורשת ממשרד החינוך שלא יקבלו את הדברים הללו – יש ללמוד את הנושאים האלו מהסיכום. באופן דומה בפרק של ''תוצאות`` מתוארות השפעות. מה שכתוב בספר על השפעות צבאיות נכונות, לא למדנו על זה. 
	
	גלבר תוקף את ההנחה שהסכמי השלום שהוצעו באמצעות ביארינג ב־71 והשנייה באמצעות קיסינג'ר ב־73 שהיו זהות להסכם של 79. הוא טוען שמדובר בשלום מומצא שלא הוצע על ידי מצרים לישראל או לארה''ב ולא נדחה על־ידן. 
	
	חשוב להבהיר שככל שיחשפו יותר מסמכים, בצד המצרי ובצד הישראלי, ככל הנראה המחקר והתפיסה של המלחמה תשתנה. 
	
	\subsubsection*{לקריאה/צפייה נוספת}
	
	\begin{itemize}
		\item \textit{שתיקת הצופרים} היא סדרה ''ישנה־ישנהההה באיכות פח`` בלטענת המורה מראה באופן יפה את מה שקדם למלחמה. 
%		\item \textit{תקומה} – 
		\item \textit{רעידת אדמה} – סדרה יהודית על מלחמת יום כיפור [שרית מתעלמת מההזהרה לתוכן גרפי ומציגה את ההתחלה של הסרטון בכיתה]
		\item \textit{שעת נעילה} – סדרת טלווזיה שמבוססת על אירועים אמיתיים. יש לשים לב שזה לא היסטוריה אלא סדרת טלווזיה. 
	\end{itemize}
	
	שימו לב שגם סדרות דוקומנטריות, בהן בוחרים מה להכניס לשם ומה לא. 
	
	[שרית משמיעה ''בוקר טובבבבב בוקר טובב`` מעוץ־לי־גוץ־לי ליפתח שנרדם]
	
	[שרית משמיעה ''לו יהי`` בביצוע של זהו זה! עם הגששים]
	
	[שרית משמיעה את ''לבחור נכון`` של אמיר דדון]
	
	נשאר לנו ללמוד עליה וקליטה, נושא ארוך וחשוב ששרית לא תלמד (זה מה שהיא אמרה, בכיתה). בין המתכונת לבגרות שרית תעשה השלמה למלחמת העצמאות. 
	
	למדנו על כמה דברים: 
	\begin{multicols}{2}
		\begin{itemize}
			\item גורמים (סיבות)
			\item נסיבות
			\item סיבות להפתעה
			\item תוצאות (טווח מידי)
			\item השפעות, שמחולקות ל־: 
			\begin{itemize}
				\item פוליטיות
				\item חברתיות
				\item כלכליות
				\item יחב''ל
				\item מדינות ערב
			\end{itemize}
		\end{itemize}
	\end{multicols}
	
	\subsection*{השפעות על מדינות ערב ששרית שכחה לדבר עליהם}
	סוריה מעולם לא רצתה מגעים לשלום. אומנם שלום לא רלוונטי, אך סוריה הייתה מארחת של ארגוני טרור לא מעטים שפעלו משטחה. מסיום המלחמה, סוריה לא יוזמת באופן ישיר מהלכים צבאיים וכביכול שומרת על הסכמי שביתת הנשק, ובאותו הזמן מאפשרת לארגוני טרור לשבת בשטחה ולהוציא פעולות טרור כנגד מדינת ישראל. 
	
	מול מצריים, לאחר הסכמי ביניים, מתחיל תהליך איטי (לאחר עליית הליכוד לשלטון), התהליך נמשך. הימין הוא זה שחתם על הסכמי שלום, בגלל שההתנגדות הייתה מהימין – השמאל בעד, לכאורה רק צריך לשכנע את הימין. 
	
	\section*{על המבחן}
	
	יכולים לשאול, לדוגמה, מה היו הגורמים למלחמה (שימו לב, הם קושרים למדינות ערב), מה היו מטרותיהן ואילו מטרות הם השיגו. לדוגמה, מצריים השיגה את הסכם השלום שהיא רצתה, מצריים השיגה את השטחים ממלחמת ששת הימים, כל מדינות ערב במלחמה (מצריים וסוריה) הצליחו להצליח להשיב את הכבוד. על אף שהן לא נצחו, ההפתעה מבחינתן היה הישג שהוביל לכך שהשיגו את הכבוד שלהם. 
	
	
	השאלות דורשות חשיבה, בגלל שיש ספר פתוח. זה הולך והופך להיות (בניסוחן) להיות פחות ראש בקיר. בתוך החומר ששרית למדה, נוכל לענות על הכל. העניין הוא להבין מה רוצים מאיתנו. לא להבהל מהניסוחין. כשמתרגלים, מחוץ לקופסה. אפשר לתרגל שאלות ''בימי אנו־בנו`` אך יש לזכור שאלו שאלות בלי ספר פתוח. כאשר יש ספר פתוח השאלות מכיוונות להבנה, לא להעתקה מהספר. 
	
	יש שאלות מכל מיני סוגים: 
	\begin{multicols}{2}
		\begin{itemize}
			\item ניתוח קטע מקור, בהן אפשר לשאול: 
			\begin{itemize}
				\item עובדה – דעה
				\item טענה (טיעון ונימוק) [לפעמים עם משפט ולא עם מקור]
				\item מידע
			\end{itemize}
			\item ניתוח שני קטעי מקור:
			\begin{itemize}
				\item השוואה [לעיתים ללא מקור]
				\item הערכת מקור
			\end{itemize}
		\end{itemize}
	\end{multicols}
	פרט לשאול מידע, הדברים יכולים להופיע בגל מני דרכים. ברוב המקרים, הטענה בשאלות טענה לא תהיה כזו שאתם מצפים לה. ''תצטרכו טיפה לשנות את הראש``. טענה היא תמיד עם נימוק. 
	
	חשוב מאוד לבנות לעצמכם את התשובה לפני שאתם כותבים אותה. ''רובכם לא עושים את זה``. כדאי לתכנן מראש מה אתם רוצים לכתוב. 
	
	בשאלת עמדה, יש להציג את עמדתכם, ולהוסיף נימוקים (עובדות היסטוריות), אם או בלי מקור. לא לשכוח לכתוב את העמדה. 
	
	בהשוואה צריך לעשות טבלת עזר, כדי לא לפספס. כיצד X התמודד עם קריטריון א' וב', כיצד Y התמודד עם קריטריון א' וב', בכל אחד מהקריטריונים יש שונה ודומה. סה''כ 6 תאים (כי דומה תא משותף בין X ל־Y). כשעושים טבלה מאוד ברור שחסר משהו, כי זה פשוט לא נמצא בטבלה. כל מי שענה בכיתה על שאלה 4, טעה בכך – לא עשו חצי מהדברים בגלל זה. לפעמים תגלו שאתם לא יודעים לעשות את ההשוואה, אז אל תעשו. 
	
	בבגרות, מותר לענות בטבלה. לא יפסלו את השאלה. ''עוד לא היה התלמיד שהצליח בתוך טבלה לכתוב תשובה מלאה`` – היא פסיכולוגית מצמצמת לכתוב את זה מצומצם בנקודות. לכן מומלץ מאוד להשתמש בטבלה כטבלת עזר, ''אלא אם כן זו טבלה שהיא עמוד``. 
	
	בבגרות עדיין יש מתווה קורונה. טוטאליטריות ושואה יש שאלת חובה, ואז שתי שאלות נוספות משני פרקים שונים. נא לא להמר. לא כדאי לא ללמוד פרק. ''שאשא ביטון, עשתה רפורמה שהייתה ונגמרה חודשיים כי היו בחירות`` – הרבה בתי ספר הורידו לאומיות, בגלל שלמדו תחת הרפורמה של שאשא ביטון. זה היה הפרק הכי קל בבגרות באותה השנה. כל שאלה שני סעיפים, עמוסים. בשנתיים האחרונות יש להם קוקווים שמבהירים מה צריך לעשות בחלוקה למטלות, במקום פסקה ארוכה. 
	
		
	
\end{document}